\section{Discussion}
\label{sec:discussion}

\subsection{Why we believe the ceiling is structural}

We have shown that more than sixty method variants on frozen
SigLIP-L features all plateau at $\kappa_q \approx 0.24$ on
DAiSEE. The phenomenon is robust to:
\begin{itemize}\itemsep0pt
  \item Encoder scale (B/14 to SO400M-14): the ceiling shifts
  upward modestly but plateaus.
  \item Encoder objective (contrastive softmax in CLIP, sigmoid in
  SigLIP, MIM in DINOv2, patch-text-alignment in TIPSv2).
  \item Probe family (linear, ordinal, non-linear MLP, ensemble,
  adapter).
  \item Pre/post-hoc calibration.
  \item Personalization (per-subject baseline, per-subject anchor
  modulation).
\end{itemize}

The mechanism is clear from the IDEP curve (Sec.~\ref{sec:idep}):
identity directions in the frozen feature space are high-rank
redundant, while engagement variance is concentrated in only
two-to-five directions that overlap with identity. Any linear
operator that removes identity also removes engagement. The
within-subject Spearman finding (Sec.~\ref{sec:within}) is the
behavioural fingerprint of a model that is reading off subject
priors rather than per-clip variation.

\textsc{Dream}, as a deliberately gentle linear identity
intervention, fails for an explainable reason: even the gentlest
identity removal disturbs the small engagement subspace. This
strengthens the structural diagnosis: the entanglement is
irreducible \emph{in the linear-readout regime}. To break the
ceiling, encoder weights must be adapted to expose
identity-orthogonal engagement directions, or non-visual signals
must be incorporated.

\subsection{What this implies for practitioners}

\paragraph{Don't expect zero-shot to work.} Across CLIP, LLaVA-1.5-7B,
and GPT-4o under three prompt variants each, zero-shot prompting
yielded $\kappa_q \in [0.00, 0.10]$. The best zero-shot is worse
than even unadjusted linear-probe baselines. Practitioners should
plan for at least a labelled subject-disjoint validation split.

\paragraph{Use the threshold-tuning recipe.} Sec.~\ref{sec:recipe}'s
recipe achieves $\kappa_q = 0.238$ from SigLIP-L features alone,
without any fine-tuning. This is the strongest frozen-feature
recipe we are aware of for DAiSEE engagement; future
encoder-fine-tuning papers should report against it.

\paragraph{Anticipate identity confound.} Whatever method is used
on a DAiSEE-style benchmark, subject identity will be near-perfectly
linearly recoverable from features. Methods that claim to be
``subject-invariant'' should report id-acc on residual features as
a sanity check, not just engagement metrics.

\paragraph{Within-subject ranking is the harder task.} Global
$\kappa_q$ is a misleading summary of how useful an engagement
classifier is for a teacher: a classifier that ranks students by
``typical engagement level'' but fails to rank a single student's
clips is of limited operational value. We recommend reporting
within-subject Spearman $\rho$ in addition to $\kappa_q$ for any
new method on DAiSEE.

\subsection{Limitations and threats to validity}

\paragraph{Dataset scope.} Our claims are scoped to DAiSEE.
Distinct engagement-recognition datasets exist
(EngageNet~\cite{khare2024engagenet}, EmotiW~\cite{dhall2017emotiw},
SCB~\cite{wang2024scb}), and the ceiling phenomenon may or may not
generalize. DAiSEE-specific factors (long uniform background, small
subject pool, label noise---see \cite{kaur2021label} for prior
critiques of the labels) may amplify the entanglement we observe.

\paragraph{Encoder choice.} We focus on contrastive
vision-language encoders. Video-pre-trained encoders
(VideoMAE-V2~\cite{wang2023videomae},
InternVideo2~\cite{wang2024internvideo2}) might encode short
temporal dynamics that frozen image encoders miss. We do not test
these in this paper; we discuss them as future work below.

\paragraph{Failure of DREAM may not generalize.} Anchor-based
personalization may help in other settings (e.g., classroom video
with controlled neutral-baseline clips), even if it fails here.
\textsc{Dream}'s failure on DAiSEE specifically should be read as
``frozen-feature anchor modulation cannot break this dataset's
ceiling,'' not as a global statement.

\paragraph{No audio.} DAiSEE is video-only; audio could carry
substantial engagement signal but is not available. This is an
important limitation we do not address.

\subsection{Future work}

\paragraph{Encoder fine-tuning with identity-adversarial
regularization.} The structural diagnosis suggests fine-tuning
\emph{the encoder} (LoRA or full last-block) on engagement labels
with a gradient-reversal subject-ID auxiliary loss should expose
identity-orthogonal engagement directions. Our preliminary LoRA
attempt diverged at epoch~1; properly stabilized
encoder-fine-tuning (warmup, gradient accumulation, longer
schedule) is the most promising direction.

\paragraph{Temporal video encoders.} The supervised SOTA
(ViBED-Net) uses 3-D temporal CNNs and reaches 73\% accuracy.
Our probe of VideoMAE-base (Kinetics pretraining, 16 frames per
clip) achieves only $\kappa_q = 0.101$---worse than image VLMs.
This suggests that frozen video features from action-recognition
pretraining do not transfer to engagement; encoders pre-trained on
affect-relevant data (FER+temporal) or supervised on engagement
itself are likely necessary. Larger video models (VideoMAE-Large,
InternVideo2) and engagement-supervised video pre-training remain
open directions.

\paragraph{Explicit behavioural fusion at the encoder level.} Our
MediaPipe-explicit signals (52-d blendshapes, 6-d gaze, 16-d body
pose) carry engagement signal at the linear-probe level but did
not exceed the frozen-feature ceiling once fused (Sec.~\ref{sec:abl-recipe}
analogue in supplementary). Building these signals into the
encoder via dual-branch training rather than fusing at the
readout level is a natural next step.

\paragraph{Audio--video multi-modal.} EngageNet provides audio;
recording audio in DAiSEE-style classroom video at deployment time
is straightforward. This is the single largest unexplored
information channel.
